% =====================================================================
\section{Neuron and synapse}
\fig{fig_neuron.png}{Neuron and synapse. On the left the whole cell, on the right the point of contact between two cells, enlarged.}
A neuron is a cell adapted to transmitting signals. It maintains an electrical voltage across its membrane. When the summed input exceeds a threshold, a short impulse runs along the output process. Between cells the signal is passed chemically: the impulse releases molecules of a \emph{transmitter} from the terminal, and these open receptors on the next cell (fig.~3).
\facts
\begin{itemize}
\item A neuron consists of a body, branching \emph{dendrites} (input) and a single \emph{axon} (output). The impulse is all-or-none; information is carried by the rate and the timing of impulses.\tU
\item The main excitatory transmitter of the cortex is \emph{glutamate}, the main inhibitory one \emph{GABA}. Excitation brings the cell closer to threshold, inhibition takes it further away.\tU
\item Glutamate has two principal receptors. \emph{AMPA} is fast, acting in milliseconds. \emph{NMDA} is slow (tens to hundreds of milliseconds); section~7 is about it.\tU
\item Calcium entering the cell during its work serves as a signal to switch genes on. In this way synaptic activity governs which proteins the cell produces.\tU
\end{itemize}
\begin{ineq}
A neuron is a threshold element, a synapse a connection with a sign and a weight. What matters more for what follows is something else: through calcium, fast electrical activity (milliseconds) governs the cell's slow chemistry (hours). This is the bridge between time scales.
\end{ineq}
